\section{Architettura hardware e software}

\subsection{Componenti software}
\begin{table}[H]
    \centering
    \begin{tabularx}{\textwidth}{@{}llX@{}}
        \toprule
        \textbf{Strato} & \textbf{Componente} & \textbf{Ruolo} \\
        \midrule
        Server & Python 3.10+ & Runtime del backend \\
        Server & Flask 3.1 + SQLAlchemy 3.1 & Framework REST e ORM \\
        Server & Flask-JWT-Extended 4.7 & Emissione/verifica dei token JWT \\
        Server & Flask-CORS 6.0 & Gestione delle policy CORS \\
        Server & SQLite (file \code{.db}) & Database relazionale embedded \\
        Server & werkzeug.security & Hashing password (PBKDF2-SHA256) \\
        Client & React Native 0.81 + Expo SDK 54 & UI mobile cross-platform \\
        Client & expo-router 6 & Routing file-based \\
        Client & expo-secure-store (opz.) & Keychain/Keystore per i token \\
        \bottomrule
    \end{tabularx}
    \caption{Stack software.}
\end{table}

\subsection{Infrastruttura}
In ambiente di sviluppo il backend gira come singolo processo Python sulla
macchina del team (port 5000). La mobile app viene servita da Metro
bundler di Expo e si collega al backend via IP della LAN.

In ambiente di produzione si assume il seguente schema di deployment
tipico:

\begin{figure}[H]
    \centering
    % \includegraphics[width=0.85\textwidth]{architettura-deployment.png}
    \fbox{\parbox{0.85\textwidth}{\centering\vspace{1em}
        \TODO{Inserire diagramma di deployment (Internet $\to$ Reverse
        Proxy con TLS $\to$ Flask/Gunicorn $\to$ SQLite/PostgreSQL).}
        \vspace{1em}}}
    \caption{Schema di deployment in produzione.}
    \label{fig:deploy}
\end{figure}
